% ============================================================
%  APPENDIX
% ============================================================
\appendix

% ------------------------------------------------------------
\section{Dataset Details}
\label{app:datasets}

Table~\ref{tab:datasets} summarises the six benchmark datasets used in this study, and Table~\ref{tab:sample_queries} provides representative sample queries from each.
All datasets were sampled to a maximum of 50 queries each, yielding 291 total
query--context--answer triples after filtering records with empty context fields.
Query-type labels were assigned automatically: HotpotQA and SQuAD v2 used
metadata flags (\texttt{type} and \texttt{is\_impossible}) embedded in the
dataset; all remaining datasets contain only single-hop questions as defined by
their original design, confirmed by manual inspection of a random 20-item subset.

\begin{table}[h]
\centering
\caption{Benchmark datasets included in the balanced evaluation corpus.}
\label{tab:datasets}
\begin{tabular}{llllr}
\toprule
\textbf{Dataset} & \textbf{Domain} & \textbf{Query Types} & \textbf{Source (HF path)} & \textbf{\textit{N}} \\
\midrule
HotpotQA           & General & Multi-hop, Comparative   & \texttt{hotpot\_qa/distractor}              & 50 \\
SQuAD v2           & General & Single-hop, Unanswerable & \texttt{rajpurkar/SQuAD-v2.0}               & 50 \\
PubMedQA (labeled) & Medical & Single-hop               & \texttt{qiaojin/PubMedQA/pqa\_labeled}      & 50 \\
PubMedQA (artif.)  & Medical & Single-hop               & \texttt{qiaojin/PubMedQA/pqa\_artificial}   & 41$^{*}$ \\
FinQA              & Finance & Single-hop               & \texttt{ibm/finqa}                          & 50 \\
FinanceBench       & Finance & Single-hop               & \texttt{PatronusAI/financebench}            & 50 \\
\midrule
\multicolumn{4}{r}{\textbf{Total}} & \textbf{291} \\
\bottomrule
\end{tabular}

\smallskip
{\small $^{*}$Nine records excluded due to empty context after preprocessing.}
\end{table}

\begin{table}[h]
\centering
\caption{Representative sample queries drawn from the six benchmark datasets.
Each row shows one question--answer pair as presented to the RAG pipeline.
``N/A'' denotes unanswerable items for which no reference answer exists.
Medical answers are truncated for brevity; full long-form answers were used
during evaluation.}
\label{tab:sample_queries}
\small
\begin{tabular}{p{2.2cm} l p{1.8cm} p{5.5cm} p{2.8cm}}
\toprule
\textbf{Dataset} & \textbf{Domain} & \textbf{Query Type} & \textbf{Question} & \textbf{Reference Answer} \\
\midrule

HotpotQA
  & General
  & Comparative
  & Were Scott Derrickson and Ed Wood of the same nationality?
  & Yes \\[4pt]

HotpotQA
  & General
  & Multi-hop
  & What government position was held by the woman who portrayed Corliss Archer in the film \textit{Kiss and Tell}?
  & Chief of Protocol \\[4pt]

SQuAD v2
  & General
  & Single-hop
  & In what country is Normandy located?
  & France \\[4pt]

SQuAD v2
  & General
  & Unanswerable
  & Who gave their name to Normandy in the 1000s and 1100s?
  & N/A \\[4pt]

PubMedQA
  & Medical
  & Single-hop
  & Do mitochondria play a role in remodelling lace plant leaves during programmed cell death?
  & Yes --- mitochondrial dynamics were observed to progress as PCD advanced \ldots \\[4pt]

PubMedQA (artif.)
  & Medical
  & Single-hop
  & Are group 2 innate lymphoid cells (ILC2s) increased in chronic rhinosinusitis with nasal polyps or eosinophilia?
  & Yes --- ILC2s are elevated in CRSwNP and may drive nasal polyp formation \ldots \\[4pt]

FinQA
  & Finance
  & Single-hop
  & What is the average payment volume per transaction for American Express?
  & 127.40 \\[4pt]

FinanceBench
  & Finance
  & Single-hop
  & What is the FY2018 capital expenditure amount (in USD millions) for 3M?
  & \$1{,}577.00 \\

\bottomrule
\end{tabular}
\end{table}


% ------------------------------------------------------------
\section{Chunking Strategy Specifications}
\label{app:chunking}

Six chunking strategies were evaluated, combining two methods (fixed-size and
semantic) with three token-budget sizes (128, 256, 512). Table~\ref{tab:chunking}
provides full specifications. All token counts use the \texttt{cl100k\_base}
tiktoken encoding (GPT-3/4 vocabulary). Fixed strategies produce non-overlapping
windows; semantic strategies respect paragraph and sentence boundaries before
falling back to fixed chunking for units that exceed the token budget.

\begin{table}[h]
\centering
\caption{Chunking strategy specifications. ``Max tokens'' denotes the upper token
limit per chunk; fixed strategies use exact-size windows with no overlap.}
\label{tab:chunking}
\begin{tabular}{llrrl}
\toprule
\textbf{Strategy} & \textbf{Method} & \textbf{Max tokens} & \textbf{Overlap} & \textbf{Boundary logic} \\
\midrule
\texttt{fixed\_128}    & Fixed    & 128 & 0   & Hard token window \\
\texttt{fixed\_256}    & Fixed    & 256 & 0   & Hard token window \\
\texttt{fixed\_512}    & Fixed    & 512 & 0   & Hard token window \\
\texttt{semantic\_128} & Semantic & 128 & --- & Paragraph $\to$ sentence $\to$ fixed \\
\texttt{semantic\_256} & Semantic & 256 & --- & Paragraph $\to$ sentence $\to$ fixed \\
\texttt{semantic\_512} & Semantic & 512 & --- & Paragraph $\to$ sentence $\to$ fixed \\
\bottomrule
\end{tabular}
\end{table}

% ------------------------------------------------------------
\section{LLM-as-Judge Prompt}
\label{app:judge_prompt}

The following prompt was used verbatim for all LLM-as-judge evaluations.
Scores were elicited at temperature 0 using \texttt{openai/gpt-4o-mini} via
the OpenRouter API. Raw integer scores (1--5) were normalized to $[0, 1]$ by
the transformation $s = (r - 1) / 4$.

\begin{quote}
\small\ttfamily\sloppy
You are an expert evaluator assessing the correctness of a RAG system's answer.

\medskip
Task\\
Rate how correctly the RESPONSE answers the QUESTION, compared to the REFERENCE ANSWER.

\medskip
Inputs\\
- QUESTION: \{question\}\\
- REFERENCE ANSWER: \{ground\_truth\}\\
- RESPONSE: \{response\}

\medskip
Scoring Rubric (1--5)\\
5 = The response contains the correct answer with the same essential information
as the reference. Minor wording differences or additional correct elaboration
are acceptable.\\
4 = The response is correct but misses minor details from the reference, or adds
minor inaccuracies alongside the correct answer.\\
3 = Partially correct. Captures some key information but misses important
details or includes significant inaccuracies.\\
2 = Mostly incorrect. Attempts to answer but misses the main point or
contradicts the reference on key facts.\\
1 = Completely incorrect, does not answer the question, or contradicts the
reference answer entirely.

\medskip
Important\\
- Focus on SEMANTIC correctness, not lexical overlap.\\
- A verbose response that contains the correct answer should score 5.\\
- Short reference answers (e.g.\ ``yes'', ``no'', a name, a date): the response
  must contain or clearly convey that answer.\\
- If the question is unanswerable and the reference indicates this, a response
  that appropriately declines should score highly.

\medskip
Respond with ONLY valid JSON: \{"score": \textless int 1--5\textgreater, "reasoning": "\textless one sentence\textgreater"\}
\end{quote}

% ------------------------------------------------------------
\section{Query-Type Classification Prompt}
\label{app:qt_prompt}

For datasets without built-in query-type metadata, the following prompt was used
to classify each question via \texttt{openai/gpt-4o-mini} at temperature 0.

\begin{quote}
\small\ttfamily\sloppy
Classify the following question into exactly one of these query types:\\
single-hop, multi-hop, comparative, unanswerable.\\
A single-hop question requires one fact.\\
A multi-hop question requires combining facts from multiple sources.\\
A comparative question requires comparing two or more entities.\\
An unanswerable question cannot be answered from typical reference material.\\
Respond with ONLY the query type label.

\medskip
Question: \{query\}
\end{quote}

% ------------------------------------------------------------
\section{RAGAS Metric Definitions}
\label{app:ragas}

Four retrieval and generation metrics were computed using the RAGAS
framework \cite{ragas2023} alongside the LLM-as-judge score. All RAGAS
metrics are bounded $[0, 1]$ (higher = better).

\begin{description}
  \item[Faithfulness] Proportion of claims in the generated answer that are
    directly supported by the retrieved context. Computed by decomposing the
    answer into atomic claims and checking each against the context with an LLM
    verifier.
  \item[Answer Relevancy] Cosine similarity between the embedding of the
    question and embeddings of synthetic questions reverse-engineered from the
    answer. Measures whether the answer actually addresses the question asked.
  \item[Context Precision] Mean reciprocal rank of relevant context chunks
    among the top-$k$ retrieved chunks, weighted by position. Measures how
    highly the retriever ranks the most relevant passages.
  \item[Context Recall] Fraction of ground-truth answer claims that can be
    attributed to the retrieved context. Measures retrieval completeness.
\end{description}

% ------------------------------------------------------------
\section{Full ANOVA Results}
\label{app:anova}

\begin{table}[h]
\centering
\caption{Type III ANOVA Results for LLM-as-Judge Score. Main effects only; interaction terms were excluded due to empty cells arising from domain--query-type confounding (medical and finance domains contain only single-hop queries). $\eta^2_p$ = partial eta squared.}
\label{tab:anova}
\begin{tabular}{lrrrrl}
\toprule
\textbf{Effect} & \textbf{F} & \textbf{df\textsubscript{num}} & \textbf{df\textsubscript{den}} & \textbf{p} & $\boldsymbol{\eta^2_p}$ \\
\midrule
Chunking strategy & 3.63 & 5 & 6336 & .003 & .003 (negligible) \\
Query type        & 30.73 & 3 & 6336 & <.001 & .014 (small) \\
Domain            & 397.66 & 2 & 6336 & <.001 & .112 (medium) \\
\bottomrule
\end{tabular}
\end{table}


% ------------------------------------------------------------
\section{Tukey HSD Post-Hoc Tables}
\label{app:tukey}

\begin{table}[h]
\centering
\caption{Tukey HSD post-hoc pairwise comparisons by domain for LLM-as-Judge Score. Mean difference = Group A $-$ Group B (positive values indicate Group A scores higher). 95\% CI = confidence interval for mean difference. Cohen's \textit{d} computed using pooled SD. ns = not significant.}
\label{tab:tukey_domain}
\begin{tabular}{llrrrl}
\toprule
\textbf{Group A} & \textbf{Group B} & \textbf{Mean Diff} & \textbf{95\% CI} & \textbf{\textit{d}} & \textbf{\textit{p}$_{\text{adj}}$} \\
\midrule
General & Finance & $+$0.257 & [$+$0.234, $+$0.281] & 0.689 & $<$.001 \\
Medical & Finance & $+$0.271 & [$+$0.245, $+$0.298] & 0.776 & $<$.001 \\
General & Medical & $-$0.014 & [$-$0.037,\hphantom{$-$}0.009] & 0.053 & .334 \\
\bottomrule
\end{tabular}
\end{table}


\begin{table}[h]
\centering
\caption{Tukey HSD post-hoc pairwise comparisons by query type for LLM-as-Judge Score. Mean difference = Group A $-$ Group B (positive values indicate Group A scores higher). 95\% CI = confidence interval for mean difference. Cohen's \textit{d} computed using pooled SD. ns = not significant.}
\label{tab:tukey_querytype}
\begin{tabular}{llrrrl}
\toprule
\textbf{Group A} & \textbf{Group B} & \textbf{Mean Diff} & \textbf{95\% CI} & \textbf{\textit{d}} & \textbf{\textit{p}$_{\text{adj}}$} \\
\midrule
Comparative & Multi-hop    & $+$0.071 & [$+$0.021,\hphantom{$+$}0.121] & 0.251 & .002 \\
Comparative & Single-hop   & $+$0.148 & [$+$0.103,\hphantom{$+$}0.193] & 0.471 & $<$.001 \\
Comparative & Unanswerable & $+$0.126 & [$+$0.074,\hphantom{$+$}0.178] & 0.436 & $<$.001 \\
Multi-hop   & Single-hop   & $+$0.077 & [$+$0.047,\hphantom{$+$}0.107] & 0.228 & $<$.001 \\
Multi-hop   & Unanswerable & $+$0.055 & [$+$0.015,\hphantom{$+$}0.094] & 0.175 & .002 \\
Single-hop  & Unanswerable & $+$0.022 & [$-$0.011,\hphantom{$+$}0.055] & 0.065 & .307 \\
\bottomrule
\end{tabular}
\end{table}


\begin{table}[h]
\centering
\caption{Tukey HSD post-hoc pairwise comparisons by chunking strategy for LLM-as-Judge Score. Mean difference = Group A $-$ Group B (positive values indicate Group A scores higher). 95\% CI = confidence interval for mean difference. Cohen's \textit{d} computed using pooled SD. All significant pairs involve \texttt{fixed\_128} as the lower-scoring group; effect sizes are negligible throughout ($d < 0.20$), consistent with the small ANOVA $\eta^2_p = .003$.}
\label{tab:tukey_strategy}
\begin{tabular}{llrrrl}
\toprule
\textbf{Group A} & \textbf{Group B} & \textbf{Mean Diff} & \textbf{95\% CI} & \textbf{\textit{d}} & \textbf{\textit{p}$_{\text{adj}}$} \\
\midrule
\texttt{fixed\_256}   & \texttt{fixed\_128}   & $+$0.042 & [$+$0.000,\hphantom{$+$}$+$0.084] & 0.122 & .049 \\
\texttt{fixed\_512}   & \texttt{fixed\_128}   & $+$0.042 & [$+$0.000,\hphantom{$+$}$+$0.084] & 0.122 & .049 \\
\texttt{semantic\_128} & \texttt{fixed\_128}  & $+$0.043 & [$+$0.001,\hphantom{$+$}$+$0.085] & 0.124 & .046 \\
\texttt{semantic\_256} & \texttt{fixed\_128}  & $+$0.056 & [$+$0.013,\hphantom{$+$}$+$0.099] & 0.163 & .003 \\
\texttt{semantic\_512} & \texttt{fixed\_128}  & $+$0.054 & [$+$0.011,\hphantom{$+$}$+$0.097] & 0.158 & .004 \\
\midrule
\texttt{fixed\_512}   & \texttt{fixed\_256}   & $+$0.000 & [$-$0.042,\hphantom{$-$}$+$0.042] & 0.000 & 1.000 \\
\texttt{semantic\_128} & \texttt{fixed\_256}  & $+$0.001 & [$-$0.041,\hphantom{$-$}$+$0.043] & 0.003 & 1.000 \\
\texttt{semantic\_256} & \texttt{fixed\_256}  & $+$0.014 & [$-$0.029,\hphantom{$-$}$+$0.057] & 0.041 & .936 \\
\texttt{semantic\_512} & \texttt{fixed\_256}  & $+$0.012 & [$-$0.030,\hphantom{$-$}$+$0.055] & 0.036 & .963 \\
\texttt{semantic\_128} & \texttt{fixed\_512}  & $+$0.001 & [$-$0.041,\hphantom{$-$}$+$0.043] & 0.003 & 1.000 \\
\texttt{semantic\_256} & \texttt{fixed\_512}  & $+$0.014 & [$-$0.029,\hphantom{$-$}$+$0.057] & 0.041 & .936 \\
\texttt{semantic\_512} & \texttt{fixed\_512}  & $+$0.012 & [$-$0.030,\hphantom{$-$}$+$0.055] & 0.036 & .963 \\
\texttt{semantic\_256} & \texttt{semantic\_128} & $+$0.013 & [$-$0.030,\hphantom{$-$}$+$0.057] & 0.038 & .954 \\
\texttt{semantic\_512} & \texttt{semantic\_128} & $+$0.012 & [$-$0.032,\hphantom{$-$}$+$0.055] & 0.033 & .975 \\
\texttt{semantic\_512} & \texttt{semantic\_256} & $-$0.002 & [$-$0.046,\hphantom{$-$}$+$0.042] & 0.005 & 1.000 \\
\bottomrule
\end{tabular}
\end{table}


% ------------------------------------------------------------
\section{Targeted Contrast: Finance vs.\ General Domain}
\label{app:contrast}

A Welch's \textit{t}-test was conducted on the difference score
(LLM Judge $-$ Faithfulness) to isolate retrieval-generation divergence
between the Finance and General domains. Results are reported in
Table~\ref{tab:targeted_contrast}.

\begin{table}[h]
\centering
\caption{Welch's \textit{t}-test: Finance vs.\ General on
  (LLM Judge Score $-$ Faithfulness). Negative mean difference indicates
  Finance scores are lower on this combined metric. Cohen's \textit{d}
  computed using pooled SD.}
\label{tab:targeted_contrast}
\begin{tabular}{lrrrrrr}
\toprule
 & \textbf{Finance} $M$ & \textbf{Finance} $SD$
 & \textbf{General} $M$ & \textbf{General} $SD$
 & \textbf{\textit{t}} & \textbf{\textit{p}} \\
\midrule
Judge $-$ Faithful. & $-$0.120 & 0.524 & $+$0.025 & 0.365 & $-$9.90 & $<$.001 \\
\bottomrule
\end{tabular}

\smallskip
{\small $N_{\text{Finance}} = 1600$, $N_{\text{General}} = 2992$,
Welch $df = 2445.9$, Cohen's $d = -0.322$ (small).}
\end{table}

This result indicates that in the Finance domain the LLM judge consistently
awarded lower scores relative to faithfulness compared with the General domain,
suggesting that Finance queries impose a correctness penalty beyond what
retrieval quality alone explains---likely reflecting the numerical and
domain-specific vocabulary demands of financial QA.

% ------------------------------------------------------------
\section{RAG Pipeline Overview}
\label{app:pipeline}

The end-to-end RAG evaluation pipeline proceeded as follows:

\begin{enumerate}
  \item \textbf{Document ingestion.} Each query's gold context passage was
    chunked according to one of the six strategies
    (Appendix~\ref{app:chunking}).
  \item \textbf{Embedding and indexing.} Chunks were embedded with
    \texttt{text-embedding-3-small} (OpenAI, $d = 1536$) and stored in a
    Qdrant vector collection, one collection per strategy.
  \item \textbf{Retrieval.} At inference time, the query was embedded with
    the same model and the top-$k = 5$ nearest chunks were retrieved via
    cosine similarity.
  \item \textbf{Generation.} Five LLMs (GPT-4o-mini, Claude 3.5 Sonnet,
    Gemini 2.5 Pro, Llama 3.1 70B, Mistral Large) generated answers from a
    fixed prompt template: \textit{``Based on the context below, answer this
    question: \{query\}$\backslash$n$\backslash$nContext:$\backslash$n\{context\}''}.
  \item \textbf{Evaluation.} Each (question, context, answer) triple was
    scored by the LLM-as-judge (Appendix~\ref{app:judge_prompt}) and by the
    four RAGAS metrics (Appendix~\ref{app:ragas}), both using
    \texttt{openai/gpt-4o-mini} via OpenRouter at temperature 0.
\end{enumerate}

All inference and evaluation calls were routed through the OpenRouter API.
LLM responses were generated with default temperature settings per model.
